\section{I. A division-free spectral pencil, all singular directions, and the moving-observation boundary}\label{i.-a-division-free-spectral-pencil-all-singular-directions-and-the-moving-observation-boundary}

\subsection{I.1 Original coordinates and the exact dual pencil}\label{i.1-original-coordinates-and-the-exact-dual-pencil}

\textbf{Source input.} The supplied \emph{Original-metric heat transitions and complete mixed-family control} (20 September 2026), P1--P6, fixes an endomorphism and a covariance path. Its positive one-parameter coefficient theorem is retained as prior work. We give a smaller direct determinant, a multiple-source extension, and individual singular-value certificates. The classical relation between polynomial coefficient geometry and root scales has precedents such as Akian, Gaubert, and Sharify (2017); the finite inequalities required here are proved below and are not a historical-priority claim.

Let the specified value space have dimension \(q\ge1\), original metric \(G_0>0\), original endomorphism \(M\), and actual positive normal-source covariances \(\Omega_i\succeq0\). Define \[
 G(\boldsymbol t)=(G_0^{-1}+\sum_{i=1}^d t_i\Omega_i)^{-1},\quad
 K_i=G_0^{1/2}\Omega_iG_0^{1/2},\quad
 C=I+\sum_i t_iK_i,\quad t_i\ge0,
\] \[
 M_0=G_0^{1/2}MG_0^{-1/2}.
 \tag{D1}
\] These are precisely covariances of orthogonalized source additions. Distinct, unorthogonalized physical sources cannot be assigned independent positive summands without first retaining their cross Gram. The isometry from \((E,G(\boldsymbol t))\) to Euclidean coordinates is \(S=C^{-1/2}G_0^{1/2}\). Thus \(T=SMS^{-1}=C^{-1/2}M_0C^{1/2}\) and \[
 TT^*=C^{-1/2}M_0CM_0^*C^{-1/2}.
\] For every \(a>0\), \[
 \boxed{\det(aI+M^{\dagger_{G(\boldsymbol t)}}M)
 =\frac{P(a,\boldsymbol t)}{D(\boldsymbol t)},\quad
 P=\det(aC+M_0CM_0^*),\quad D=\det C.}
 \tag{D2}
\] Indeed \(T^*T\) and \(TT^*\) have the same characteristic polynomial by \(\det(I+UV)=\det(I+VU)\), which follows by eliminating either diagonal block of a two-by-two block matrix. Multiplying the determinant of \(aI+TT^*\) by \(\det C\) proves D2. This uses no covariance eigenbasis and no adjugate division. At \(a=0\), continuity gives \[
 P(0,\boldsymbol t)=|\det M|^2D(\boldsymbol t).
 \tag{D3}
\] The source's P6 adjugate formula is equal to D2, not a different determinant. In particular its exact divisibility follows from this \(q\)-square pencil.

\subsection{I.2 Several noncommuting source additions still have positive coefficients}\label{i.2-several-noncommuting-source-additions-still-have-positive-coefficients}

Choose any factors \(K_i=B_iB_i^*\) on their actual positive ranges. Put \(B_0=I\), \(t_0=1\). Then \[
 aC+M_0CM_0^*=
 \sum_{i=0}^d t_i\{aB_iB_i^*+(M_0B_i)(M_0B_i)^*\}.
 \tag{D4}
\] Expand the Gram determinant by Cauchy--Binet, using the fixed columns \(B_i,M_0B_i\) and scalar column weights \(at_i,t_i\). Every coefficient is a sum of squared moduli of \textbf{complete} \(q\)-column minors; all phases inside those determinants are retained. Therefore \[
 P(a,\boldsymbol t)=\sum_{j,\alpha}p_{j,\alpha}a^j\boldsymbol t^\alpha,
 \quad p_{j,\alpha}\ge0,\quad 0\le j\le q,\quad |\alpha|\le q.
 \tag{D5}
\] Also \(\alpha_i\le\min(q,2\operatorname{rank}K_i)\), because the two groups from source \(i\) have at most \(2\operatorname{rank}K_i\) columns. No commutativity of the \(K_i\) is used. The coefficient field is that of the original matrices and their conjugates: the square-root factors above prove positivity but need not be computed to form D2.

At most \[
 N_{q,d}=(q+1)\binom{q+d}{d}
\] coefficients occur. On positive parameters, if \(L\) is the maximum logarithm of their evaluated nonzero monomials, \[
 0\le\log\det(aI+M^\dagger M)-[L-\log D]\le\log N_{q,d}.
 \tag{D6}
\] At four equal-dimensional endpoints with the same parameters and signs \(+,+,-,-\), the absolute signed error is at most \(2\log N_{q,d}\). For fixed \(d\) this is \(O_d(\log(q+1))\). It is not dimension-independent if the number of separately penalized sources grows.

\subsection{I.3 Recover each squared singular value, including the smallest}\label{i.3-recover-each-squared-singular-value-including-the-smallest}

Let \(\lambda_1\ge\cdots\ge\lambda_q\ge0\) be the eigenvalues of \(M^\dagger M\) in the actual metric, and write \[
 P(a,\boldsymbol t)=\sum_{r=0}^q a^{q-r}A_r(\boldsymbol t),\qquad
 e_r(\lambda)=A_r/D,\quad e_0=1.
\] Let \(R=\operatorname{rank}M\). Then \(A_r>0\) for \(0\le r\le R\) and \(A_r=0\) for \(r>R\), throughout the nonnegative parameter domain. For \(1\le r\le R\), \[
 \boxed{\frac{e_r}{(q-r+1)e_{r-1}}
 \le\lambda_r\le r\frac{e_r}{e_{r-1}}.}
 \tag{D7}
\] \textbf{Proof.} For each \((r-1)\)-subset of indices, add a missing member of \(\{1,\ldots,r\}\). The product increases by at least \(\lambda_r\), and each resulting \(r\)-subset has at most \(r\) preimages. Thus \(re_r\ge\lambda_re_{r-1}\). Conversely remove from every \(r\)-subset a member whose index is at least \(r\). Its value is at most \(\lambda_r\), and each remaining subset has at most \(q-r+1\) possible extensions. This gives \(e_r\le(q-r+1)\lambda_re_{r-1}\). Zeros cause no exception when \(r\le R\). This proves D7.

At fixed parameters let \(m_r\) be the largest evaluated nonzero monomial of \(A_r\), and let \(n_r\) count its nonzero evaluated monomials. Then \(m_r\le A_r\le n_rm_r\), and the denominator \(D\) cancels \textbf{exactly} between adjacent coefficients. Hence \[
 \boxed{\frac{m_r}{(q-r+1)n_{r-1}m_{r-1}}
 \le\lambda_r\le
 \frac{r n_rm_r}{m_{r-1}}.}
 \tag{D8}
\] Here \(n_r\le\binom{q+d}{d}\). In particular, for one activation variable, \[
 \left|\log\lambda_r-[\log m_r-\log m_{r-1}]\right|
 \le\log(q(q+1)).
 \tag{D9}
\] For squared-singular logarithms divided by \(2q\), the uncertainty in D9 is \(O(\log q/q)\), for \textbf{every rank}, including the smallest positive one. This is a finite bound; no convergence or coefficient asymptotic has been inferred.

The constants in D7 cannot uniformly be replaced by one. For equal first \(r\) roots and vanishing later roots, its upper factor \(r\) is attained. If the first \(r-1\) roots tend to infinity relative to the remaining equal roots, its lower factor \(q-r+1\) is approached.

\subsection{I.4 Heat, inverse exterior powers, and simultaneous scales}\label{i.4-heat-inverse-exterior-powers-and-simultaneous-scales}

Let \(l_r,u_r\) be the bounds in D8. For \(\tau\ge0\), \[
 \sum_{r\le R}(1-e^{-\tau l_r})
 \le q-\operatorname{Tr}e^{-\tau M^\dagger M}
 \le\sum_{r\le R}(1-e^{-\tau u_r}),
 \tag{D10}
\] \[
 \sum_{r\le R}\log(1+\tau l_r)
 \le\log\det(I+\tau M^\dagger M)
 \le\sum_{r\le R}\log(1+\tau u_r).
 \tag{D11}
\] Zeros contribute zero to each deficit and logarithm, and one each to the heat trace itself. The same monotonicity gives bounds for \(\tau\lambda/(1+\tau\lambda)\). If \(M\) is invertible, there is a sharper direct return to the original inverse-exterior calculation: \[
 \boxed{\|\wedge^p M^{-1}\|_{\mathrm{HS},G}^{2}
       =\frac{A_{q-p}}{A_q},\qquad
 \frac{A_{q-p}}{\binom qp A_q}
 \le\|\wedge^pM^{-1}\|_G^2
 \le\frac{A_{q-p}}{A_q}.}
 \tag{D11a}
\] The squared singular values of the inverse exterior map are the products of \(p\) reciprocals. Their sum is \(e_p(\lambda^{-1})=e_{q-p}(\lambda)/e_q(\lambda)=A_{q-p}/A_q\). Their maximum lies between their average and their sum. This proves D11a, including \(p=0,q\). It retains the exact full root determinant through \(A_q=|\det M|^2D\). No singular direction is removed. For a separately prescribed conductor between two different metrics, the matching pullback and coordinate maps must be applied before this endomorphism formula is used.

For one activation \(t=e^{-2q\theta}\) define \[
 z_r(\theta)=\frac1{2q}
 \left[\max_n(\log p_{q-r,n}-2q\theta n)
       -\max_n(\log p_{q-r+1,n}-2q\theta n)\right].
\] This is a finite piecewise-affine function, and D9 bounds the exact logarithmic singular scale within \(\log(q(q+1))/(2q)\). At \(\tau=e^{-2q\sigma}\), these intervals certify which directions lie above or below the heat threshold. A direction whose interval meets \(\sigma\) remains unresolved. The actual coefficients are needed; the known covariance eigenvalue law alone does not determine them.

For example, fix \(M=\left(\begin{smallmatrix}1&A\\0&1\end{smallmatrix}\right)\) and let \(D=1+t\kappa\). The two covariances \(C_+=\operatorname{diag}(D,1)\) and \(C_-=\operatorname{diag}(1,D)\) have the same eigenvalues and determinant, but \[
 P_+=(a+1)^2D+aA^2,\qquad
 P_-=(a+1)^2D+aA^2D^2.
 \tag{D12}
\] Their squared-singular sums are \(2+A^2/D\) and \(2+A^2D\), respectively, while both products are one. This exhibits the orientation information retained by D2.

\subsection{I.5 The moving observation needs a different pencil}\label{i.5-the-moving-observation-needs-a-different-pencil}

Use exactly the supplied M1--M4 block frame. Write \[
 C(t)=\begin{pmatrix}I+tU&tV\\tV^*&D_t\end{pmatrix},\quad D_t=I+tW,
\quad X_t=tVD_t^{-1},
\] with the old action blocks \(\left(\begin{smallmatrix}a_0&b_0\\c&d_0\end{smallmatrix}\right)\). The actual observed endomorphism and metric are \[
 B_t=d_0+cX_t=Y_tD_t^{-1},\quad
 Y_t=d_0D_t+t cV,\quad Q_t=D_t^{-1}.
\] D2 applied at each fixed \(t\) proves \[
 \boxed{\det(aI+B_t^{\dagger_{Q_t}}B_t)
 =\frac{\det\begin{pmatrix}D_t&Y_t^*\\-Y_t&aD_t\end{pmatrix}}
              {\det(D_t)^2}.}
 \tag{D13}
\] The block determinant follows from its Schur complement. It is polynomial of degree at most \(r\) in \(a\) and \(2r\) in \(t\), where \(r\) is the observed dimension. At a fixed real \(t\ge0\) its \(a\)-coefficients are nonnegative. Its \textbf{joint} coefficients in \((a,t)\) need not be nonnegative, because \(Y_t\) contains the moving-section phases.

Take \(K=\left(\begin{smallmatrix}1&1\\1&1\end{smallmatrix}\right)\), \(M=\left(\begin{smallmatrix}2&1\\-3&1\end{smallmatrix}\right)\), and observe the second coordinate. Then \(D_t=1+t\), \(Y_t=1-2t\), and the numerator in D13 is \[
 a(1+t)^2+(1-2t)^2.
 \tag{D14}
\] Its coefficient of \(a^0t^1\) is \(-4\). The full \(M\) has determinant five, yet its observed endomorphism vanishes at \(t=1/2\). Thus neither full-rank return nor positive joint coefficients transfer by fiat to the moving observation.

Even the full scalar polynomial does not determine a fixed projected current. With \(K=\operatorname{diag}(1,2)\), the same observation, \(x=(1,1)^T\), and \[
 A_+=\begin{pmatrix}0&1\\1&0\end{pmatrix},\qquad
 A_-=\begin{pmatrix}0&i\\-i&0\end{pmatrix},
\] the full spectral polynomials, covariance, and eigenvalues agree, while the source C1 current \(\Re(\bar z w)\) is respectively \[
 \frac{1}{(1+t)(1+2t)},\qquad -\frac{1}{(1+t)(1+2t)}.
 \tag{D15}
\] The two actions are related by a diagonal unitary commuting with the covariance; transporting the class by that unitary would preserve the current. Here the class is fixed, exactly exposing the extra data lost by the scalar spectral invariant. These are auxiliary finite examples, not xi-packet evaluations.
